\cotefile{136}
\scanpage{297}
\[
\begin{aligned}
&\deleted{=\bigl(\underline{\mathrm{Hom}}^{*}((\tau_N\unsure{\mathcal{B}|B|S}))[-i],Q)\bigr)_0}\\
&\deleted{=\underline{\mathrm{Hom}}^{*}(\tau_{N-i}S)}\\
&\simeq\bigl(\underline{\mathrm{Hom}}^{*}(\tau_{N-i}(\mathcal{B}),\mathcal{A})\bigr)_i\\
&\simeq\underline{\mathrm{Hom}}^{i}(\tau_{N-i}(\mathcal{B}),\mathcal{A})
\end{aligned}
\]

Enfin si dans $C$ les $\underline{\lim}$ existent, alors $i_N$ admet un adjoint à gauche $q_N$.

En dualisant (appliquant ce qui précède à $C^o$, en utilisant l'isom
\diagram{\begin{tikzcd}[ampersand replacement=\&,column sep=small]
\mathcal{M}(C,S)^o \arrow[dd] \& \simeq \& \mathcal{M}(C^o,S)\arrow[dd] \\
\& \& \\
\mathcal{M}^{N}_{\deleted{\unsure{N|\leq N}}}(C,S)^o \& \simeq \& \mathcal{M}_{\leq N}(C^o,S)
\end{tikzcd}}
on voit que (si $C$ \unsure{ou} $C^o$ \unsure{est} une cat. de Groth. p.ex) le foncteur
\[
i^N:\mathcal{M}^N\longrightarrow\mathcal{M}
\]
\deleted{admet} \inserted{s'insère dans} une suite \deleted{\unsure{d'adjoints}} \inserted{de 4} foncteurs adjoints
\[
\varphi^{N}\ \tau^{N}\ i^{N}\ q^{N}
\]
qui se décrivent de façon duale \unsure{en|---}

% archive sheet number: 298
